% from ega3-6-fr.tex
\oldpage[III]{151}
$\mathscr O_X$-Modules (resp. de $\mathscr O_Y$-Modules) quasi-cohérents limité
inférieurement. Supposons que tous les $\mathscr P_i$ soient $S$-plats. Alors
il existe un isomorphisme canonique de $\partial$-foncteurs en
$\mathscr Q_\bullet$
\[
  \mathscr{T}\!or_\bullet^S(\mathscr P_\bullet,\mathscr Q_\bullet)
  \xrightarrow{\sim}
  \mathscr H_\bullet(\mathscr P_\bullet\otimes_S\mathscr Q_\bullet).
  \tag{6.5.9.1}\label{III.6.5.9.1.fr}
\]
\end{namedtheorem}

\begin{proof}
Ce n'est autre que (6.3.7) lorsque $S$, $X$, $Y$ sont affines; on passe de là
au cas général par les raisonnements de (6.5.2) et (6.5.3).
\end{proof}

\begin{namedtheorem}[6.5.10]{Corollaire}
Supposons que $X$ soit plat sur $S$ ($0_{\mathrm I}$, 6.7.1), que
$\mathscr P_\bullet$ et $\mathscr Q_\bullet$ soient limités inférieurement et
que tous les $\mathscr P_i$ soient des $\mathscr O_X$-Modules localement
libres (non nécessairement de type fini). Alors l'homomorphisme (6.5.9.1) est
bijectif.
\end{namedtheorem}

\begin{proof}
En effet, l'hypothèse de (6.5.9) est remplie, la platitude étant une propriété
ponctuelle sur $X$ par définition et toute somme directe de modules plats étant
un module plat ($0_{\mathrm I}$, 6.1.2).
\end{proof}

\begin{proposition}[6.5.11]
Soient $X'$, $Y'$ deux $S$-préschémas, $f:X\to X'$, $g:Y\to Y'$ deux
$S$-morphismes affines. Soit $\mathscr P_\bullet$ (resp.
$\mathscr Q_\bullet$) un complexe de $\mathscr O_X$-Modules (resp. de
$\mathscr O_Y$-Modules) quasi-cohérents; on a alors un isomorphisme canonique
fonctoriel
\[
  (f\times_S g)_*
  \bigl(\mathscr{T}\!or_\bullet^S
    (\mathscr P_\bullet,\mathscr Q_\bullet)\bigr)
  \xrightarrow{\sim}
  \mathscr{T}\!or_\bullet^S
  (f_*(\mathscr P_\bullet),g_*(\mathscr Q_\bullet)).
  \tag{6.5.11.1}\label{III.6.5.11.1.fr}
\]
\end{proposition}

\begin{proof}
Comme $f$ et $g$ sont affines, $f_*(\mathscr P_\bullet)$ et
$g_*(\mathscr Q_\bullet)$ sont des complexes de Modules quasi-cohérents
(II, 1.2.6), et si l'on pose $Z'=X'\times_S Y'$, les deux membres de (6.5.11.1)
sont des $\mathscr O_{Z'}$-Modules quasi-cohérents (6.5.1); on se ramène
aisément au cas où $S$, $X'$ et $Y'$ sont affines; mais alors il en est de même
par hypothèse de $X$ et $Y$ et la vérification résulte aussitôt de (6.4.1.1) et
(I, 1.6.3).
\end{proof}

\begin{namedtheorem}[6.5.12]{Remarque}
Soient $X'$, $Y'$ deux $S$-préschémas et supposons que $X'$ soit $S$-plat;
soient $X$ un sous-préschéma fermé de $X'$, $i:X\to X'$ l'injection canonique,
$\mathscr P_\bullet$ (resp. $\mathscr Q_\bullet$) un complexe de
$\mathscr O_X$-Modules (resp. de $\mathscr O_{Y'}$-Modules) quasi-cohérents,
limité inférieurement. Soit enfin $\mathscr L'_{\bullet\bullet}$ une résolution
de $i_*(\mathscr P_\bullet)$ formé de $\mathscr O_{X'}$-Modules localement
libres, telle que tout point de $X'$ ait un voisinage ouvert affine $U$ pour
lequel $\mathscr L'_{j,\bullet}|U$ soit une résolution libre de
$i_*(\mathscr P_j)|U$ pour tout $j$. On a alors un isomorphisme canonique
\[
  (i\times_S 1)_*
  \bigl(\mathscr{T}\!or_\bullet^S
    (\mathscr P_\bullet,\mathscr Q_\bullet)\bigr)
  \xrightarrow{\sim}
  \mathscr H_\bullet
  (\mathscr L'_{\bullet\bullet}\otimes_S\mathscr Q_\bullet).
  \tag{6.5.12.1}\label{III.6.5.12.1.fr}
\]

Si $S$, $X'$, $Y'$ sont affines et si $\mathscr L'_{j,\bullet}$ est une
résolution libre de $i_*(\mathscr P_j)$ pour tout $j$, on est ramené, en vertu
de (6.5.11), au cas où $X'=X$ est $S$-plat, et il suffit d'appliquer (6.3.4).
Dans le cas général, on définit localement l'isomorphisme (6.5.12.1), et il
s'agit de vérifier que cette définition donne bien un isomorphisme global.
Pour cela, il faut se reporter à la définition du premier isomorphisme
(6.3.4.1) qui provient d'un isomorphisme de suites spectrales (0, 11.6.5 et
11.5.3), obtenu lui-même à partir d'un morphisme de bicomplexes
$L_{\bullet\bullet}\to L''_{\bullet\bullet}$, où $L_{\bullet\bullet}$ est la
résolution donnée de $P_\bullet$, $L''_{\bullet\bullet}$ une résolution
projective de $P_\bullet$ dans la catégorie des complexes de $A$-modules
limités

